The archived rat-hindlimb model of \citet{deng_neuromechanical_2019} provides an original-engine demonstration of the experimental workflow. The model was exported to standalone AnimatLab simulation form and instrumented with neural-voltage charts. In a \qty{10.1}{\second} run without external neural drive, rhythm-generator extensor and flexor populations alternate with an approximately \qty{0.23}{\second} period, while the two legs fire in phase. The baseline retains the archived model's mechanical torque input.

A \qty{2}{\nano\ampere} current was applied to the left rhythm-generator flexor from \SIrange{2}{2.1}{\second}, using the stimulation timing reported by \citet{deng_biarticular_2023}. A flexor burst begins at \qty{2.000}{\second}, compared with \qty{2.433}{\second} in the baseline. The post-stimulus period is \qty{0.224}{\second}, compared with \qty{0.221}{\second} in the baseline. Applying the same current and timing to the knee pattern-formation flexor produces a post-stimulus rhythm-generator phase difference of \qty{-2.3 \pm 2.4}{\percent} of the baseline period. An additional, longer knee-pattern-formation stimulus from \SIrange{2}{2.5}{\second} illustrates the motor-output response in Figure~\ref{fig:deng_original_perturbation}D. These are observations from the archived antagonist model in its original engine, rather than a quantitative validation of the Simulink translation or a reproduction of the later biarticular model's results.
\begin{figure}[p]
\centering
\includegraphics[width=\textwidth]{deng_original_reviewed.pdf}
\caption[Archived Deng-model neural perturbations]{Neural perturbations in the archived antagonist model, simulated in its original AnimatLab engine. Solid lines show the baseline; dashed lines show stimulated conditions. (A, B)~Left rhythm-generator flexor and extensor responses to rhythm-generator flexor stimulation. (C)~Left rhythm-generator extensor response to knee pattern-formation flexor stimulation. The current is \qty{2}{\nano\ampere} from \SIrange{2}{2.1}{\second} in A--C. (D)~Left knee extensor motoneuron response to an extended knee pattern-formation flexor stimulus from \SIrange{2}{2.5}{\second}. Shading marks each panel's actual stimulus window.}
\label{fig:deng_original_perturbation}
\end{figure}
